\hspace{3mm}

\centerline{\bf \Huge Тренировка 2.}

\begin{flushright}
        \scriptsize{В науке сила человечества.\\ Гендо Икари.}
\end{flushright}


\problem{1}
Назовём натуральное число \textit{интересным}, если сумма его цифр --- простое число. Какое наибольшее количество подряд идущих интересных чисел можно найти?

\problem{2}
Из чисел от $1$ до $100$ выбрано  $51$  число. Докажите, что среди выбранных чисел найдутся два, одно из которых делится на другое.

\problem{3}
Решите уравнение в целых числах $x^2 + y^2 = 2022$

\problem{4}
Найдите хотя бы 2024 решения уравнения $y^2 = x^2 + x^3$ в целых числах

\problem{5}
Выпуклый четырёхугольник $ABCD$ таков, что $AC = DC$ и $\angle DBC + \angle + ABC = 180^{\circ}$. Докажите, что что этот четырёхугольник вписанный

\problem{6}
В треугольнике $ABC$ провели биссектрису $BL$ и продлили до пересечения с описанной окружностью в точке $K$. Оказалось, что $ABL$ -- равнобедренный треугольник. Докажите, что $BCK$ -- равнобедренный треугольник.

\problem{7}
Четырехугольник $ABCD$ вписан в окружность. Известно, что $AB = BC = CD < DA$. Касательная к окружности в точке $B$ пересекается с прямой $AD$ в точке $E$. Точка $F$ диаметрально противоположна точке $C$. Докажите, что треугольник $DEF$ -- равнобедренный.

\problem{8}
Найдется ли такое простое число p, большее $10^{100}$, что число $p + 210$ --- составное?
